\documentclass[english, parskip=full]{scrartcl}
\usepackage[utf8]{inputenc}  % Kodierung der Datei
\usepackage[english]{babel}  % Sprache auf Deutsch 
\usepackage{color}
\usepackage{graphicx, gensymb}
\usepackage{amsfonts, cancel, amssymb, slashed}
\usepackage{amsmath, amsthm, array, esvect, esint}
\usepackage{booktabs, multirow}  % schönere Tabellen
\usepackage{caption}  % Anpassung der Captions
\usepackage{scrlayer-scrpage}    % Manipulation des Seitenstils
\pagestyle{scrheadings}  % Stil für die Seite setzen
\automark{section}  % Abschnittsnamen als Seitenbeschriftung 
\setheadsepline{.5pt}    % Kopzeile durch Linie abtrennen
\usepackage[hidelinks]{hyperref}  % Links und weitere PDF-Features
\usepackage{typearea, tikz, pgffor, verbatim}
\usetikzlibrary{calc, quotes}
\usepackage[cache=false, outputdir=build]{minted}
\usepackage{dsfont}


\definecolor{bg}{rgb}{0.9,0.9,0.9}
\newcommand\numberthis{\addtocounter{equation}{1}\tag{\theequation}}
\newcommand{\bra}{}
\newcommand{\kom}{}
\newcommand{\brak}{}
\newcommand{\braket}{}
\newcommand{\intx}{}
\newcommand{\intr}{}
\newcommand{\kla}{}
\newcommand{\klam}{}
\newcommand{\klammer}{}
\newcommand{\oper}{}
\newcommand{\tecxt}{}
\newcommand{\Bra}{}
\newcommand{\Ket}{}
\def\I{\text{i}}
\def\E{\text{e}}

\newcommand*{\titel}{06 Potential wells and spin}
\newcommand*{\autor}{Joachim Schwardt}

\hypersetup{pdfauthor={\autor}, pdftitle={\titel}}

\subject{Quantum theory 2}
\title{\titel}
\author{\autor}
\date{\today}

\begin{document}

%\begin{titlepage}
\maketitle  % Titel setzen
%\tableofcontents  % Inhaltsverzeichnis setzen
%\end{titlepage}

\section{Questions}
Possible operators: $x_1+x_2$, $(x_1-x_2)^2$, $p_1\cdot p_2$. 
Not possible is $x_1-x_2$ (but why?)
\\
We have $E\phi=-\frac{1}{2m}\partial_x^2\phi$, which leads to $\phi=A\E^{\I kx}+B\E^{-\I kx}$ with $k^2=2mE$.
Boundary conditions are $\phi(0)=0=\phi(L)$.
The first leads to $B=-A$, or $\phi\sim \sin kx$, and the second leads to $kL=n\pi$ for $n\in\mathbb{Z}\setminus\{0\}$.
With normalization, the solutions can be written as
\begin{align*}
\phi_n&=\sqrt{\frac{2}{L}}\sin\frac{n\pi x}{L}
\end{align*}
and have energies 
\begin{align*}
E&=\frac{k^2}{2m}=\frac{n^2\pi^2}{2mL^2}
\end{align*}
The hydrogen molecule hamiltonian without spins is
\begin{align*}
H&=H_{1a}+H_{2b}+\alpha\kla\left( {\frac{1}{r_{ab}}+\frac{1}{r_{12}}-\frac{1}{r_{1b}}-\frac{1}{r_{2a}}} \right)•,
\end{align*}
where $A=a,b$ label the proton coordinates and $j=1,2$ label the electron coordinates.
Here $H_{1a,2b}$ are the hydrogen atom hamiltonians, given by
\begin{align*}
H_{j,A}&=-\frac{1}{2m}\nabla_j^2-\frac{\alpha}{r_{j,A}}.
\end{align*}
In the Heitler-London approximation, the used are eigenstates are those of the hydrogen atom, i.e. $\psi_{nlm,A}$.
\section{Two identical particles in a potential well}
This problem can more easily be solved in the coordinate system used above.
The resulting energy levels are of course identical, and the only necessary transformation is given by $x\rightarrow x-r_0$ with $L:=2r_0$.
The Hamiltonian is simply given by a sum over the two one-particle Hamiltonians:
\begin{align*}
H&=\sum_j -\frac{1}{2m}\nabla_j^2+V(x_j)
\end{align*}
The solutions trivially factor into spin and spacial parts via $\Ket | {\psi} \rangle=\Ket | {\phi} \rangle\Ket | {\chi} \rangle$, where we are given that $\Ket | {\chi} \rangle\equiv\Ket | {Sm_S} \rangle^+$ is symmetric.
We can split the basis $\phi_n$ into symmetric and antisymmetric eigenstates by defining 
\begin{align*}
\phi_n^+&:=\sqrt{\frac{2}{L}}\sin\frac{(2n-1)\pi x}{L}&\tecxt\text{and}&&\phi_n^-&:=\sqrt{\frac{2}{L}}\sin\frac{2n\pi x}{L},
\end{align*}
which have energies 
\begin{align*}
E_n^+&=\frac{(2n-1)^2\pi^2}{2mL^2}&\tecxt\text{and}&&E_n^-&=\frac{(2n)^2\pi^2}{2mL^2}.
\end{align*}
We can construct the special two-particle wave functions by symmetrisation.
Let $n,n'$ be the energy quantum numbers and $p,p'$ be the parity quantum numbers. 
Then we can define
\begin{align*}
\Ket | {\phi} \rangle^\pm &:=\frac{1}{\sqrt{2}}\kla\left( {\phi_n^p(x_1)\phi_{n'}^{p'}(x_2)\pm \phi_n^p(x_2)\phi_{n'}^{p'}(x_1)} \right)
\end{align*}
For bosons, only the symmetric solution is allowed. 
As such, we find the groundstate for $n=n'=1$ and $p=p'=+1$, since these give the lowest one-particle energy solutions.
The energy is
\begin{align*}
E_\tecxt\text{GS,B}&=2E_1^+=\frac{\pi^2}{mL^2}.
\end{align*}
The allowed spacial wave function for fermions must be antisymmetric, considering the spin assumption.
Let us list the first few lowest one-particle energies by defining $\epsilon_n^\pm:=\frac{2mL^2E_n^\pm}{\pi^2}$, i.e. $\epsilon_n^+=(2n-1)^2$ and $\epsilon_n^-=(2n)^2$:
\begin{align*}
\epsilon_n^+&=\{1,9,25,49,\dots\}&\tecxt\text{and}&&\epsilon_n^-&=\{4,16,36,64,\dots\}.
\end{align*}
We find the groundstate for $n=1,n'=1$ and $p=+,p'=-$, which corresponds to an energy of
\begin{align*}
E_\tecxt\text{GS,F}&=E_1^+ + E_1^- =\frac{5\pi^2}{2mL^2}.
\end{align*}
\section{Hydrogen molecule without spin}
This calculation is very elaborate. 
Note that operators really only act on the space, that is indicated by their superscript.
Everything else should be straight forward.
The result is identical to the lecture, we find $c_2=\pm c_1$.
\section{Hydrogen molecule with spin}
Here it is crucial to note that $S_1S_2=\frac{1}{2}(S^2-s_1^2-S_2^2)$, and that $S^2\Ket | {Sm_S} \rangle=S(S+1)\Ket | {Sm_S} \rangle$
Thus the eigenvalues of this spin product operator are
\begin{align*}
S_1S_2\Ket | {Sm_S} \rangle^+&=S_1S_2\Ket | {1 m} \rangle=\frac{1}{2}\kla\left( {1(1+1)-\frac{1}{2}\kla\left( {\frac{1}{2}+1} \right)-\frac{1}{2}\kla\left( {\frac{1}{2}+1} \right)} \right)\Ket | {1 m} \rangle=\frac{1}{4}\Ket | {1 m} \rangle\\
S_1S_2\Ket | {Sm_S} \rangle^-&=S_1S_2\Ket | {0 0} \rangle=\frac{1}{2}\kla\left( {0-\frac{3}{4}-\frac{3}{4}} \right)\Ket | {00} \rangle=-\frac{3}{4}\Ket | {00} \rangle.
\end{align*}
Define $H_\tecxt\text{spin}:=\alpha +\beta S_1S_2$. 
Then we can determine $\alpha,\beta$ from the given constraint of 
\begin{align*}
H_\tecxt\text{spin}\Ket | {Sm_S} \rangle^\mp &\overset{!}{=}E_\pm\Ket | {Sm_S} \rangle^\mp
\end{align*}
by realizing that these spin states are eigenstates of our constructed Hamiltonian, i.e.
\begin{align*}
{}^+\braket\langle {Sm_S} | {H_\tecxt\text{spin}} | {Sm_S}\rangle^+&=\alpha + \frac{\beta}{4}=E_-\\
{}^-\braket\langle {Sm_S} | {H_\tecxt\text{spin}} | {Sm_S}\rangle^-&=\alpha - \frac{3\beta}{4}=E_+.
\end{align*}
This set of equations can be solved by inverting the corresponding coefficient matrix
\begin{align*}
\begin{pmatrix}
E_+ \\ E_-
\end{pmatrix}&=\begin{pmatrix}
1 & -3/4 \\ 1 & 1/4
\end{pmatrix}\begin{pmatrix}
\alpha \\ \beta
\end{pmatrix}&\Rightarrow &&\begin{pmatrix}
\alpha \\ \beta
\end{pmatrix}&=\begin{pmatrix}
 1/4 & 3/4 \\ -1 & 1
\end{pmatrix}\begin{pmatrix}
E_+ \\ E_-
\end{pmatrix}.
\end{align*}
This gives us
\begin{align*}
H_\tecxt\text{spin} &=\frac{E_+ +3E_-}{4}+(E_- - E_+)S_1S_2.
\end{align*}

\end{document}